% Ficha de caso de uso:
%   \fichacu{ID}{Nombre}{Actor}{Precondicion}{items del flujo}{Postcondicion}{Excepcion o nota}
\newcommand{\fichacu}[7]{%
  \subsubsection*{#1. #2}
  \noindent\textbf{Actor:} #3\par
  \noindent\textbf{Precondición:} #4\par
  \noindent\textbf{Flujo principal:}
  \begin{enumerate}[leftmargin=2.2em,itemsep=0.1em,topsep=0.2em]
  #5
  \end{enumerate}
  \noindent\textbf{Postcondición:} #6\par
  \if\relax\detokenize{#7}\relax\else\noindent\textbf{Excepción o nota:} #7\par\fi
  \medskip}

\section{Especificación de Requerimientos}
\label{sec:requerimientos}

Este apartado describe el prototipo tal como está implementado, no una versión aspiracional.
Adopta la numeración de la hoja maestra de requerimientos del equipo: dieciocho casos de uso
(CU-01 a CU-18), dieciocho requisitos funcionales (RF-01 a RF-18) y ocho requisitos no
funcionales (RNF-01 a RNF-08), relacionados en la matriz de trazabilidad del
apartado~\ref{subsec:trazabilidad}. El estado de cada requisito se contrastó con el código del
prototipo vigente. Donde la hoja maestra describía una versión anterior del sistema (el detector
seleccionable en CU-03 y RF-03, el desempate por color en CU-04, la cifra de inferencia de
RNF-06 y el estado de RNF-01), este informe recoge el comportamiento actual. La
reidentificación por apariencia y la asociación entre cámaras se documentan en CU-04 y RF-04.

\subsection{Requerimientos funcionales}
\label{subsec:rf}

\subsubsection{Casos de uso del sistema}
\label{subsec:casos-uso-tabla}

La Tabla~\ref{tab:casos-uso} resume los dieciocho casos de uso y la
Figura~\ref{fig:casos-uso} los relaciona con sus actores. El caso CU-16 conserva su número en
la hoja maestra, pero corresponde a un módulo retirado del prototipo vigente.

{\small
\setlength{\LTleft}{0pt}\setlength{\LTright}{0pt}
\begin{longtable}{@{}L{1.3cm}L{6.0cm}L{3.8cm}L{2.6cm}@{}}
\caption{Casos de uso del sistema AeroTrack.}\label{tab:casos-uso}\\
\toprule
\textbf{ID} & \textbf{Caso de uso} & \textbf{Actor} & \textbf{Estado} \\
\midrule
\endfirsthead
\toprule
\textbf{ID} & \textbf{Caso de uso} & \textbf{Actor} & \textbf{Estado} \\
\midrule
\endhead
\bottomrule
\endlastfoot
CU-01 & Configurar el plano y las cámaras & Operador & Implementado \\
CU-02 & Calibrar una cámara & Operador & Implementado \\
CU-03 & Procesar una fuente y hacer seguimiento & Operador & Implementado \\
CU-04 & Asociar identidad entre cámaras y reidentificar personas & Sistema & Implementado \\
CU-05 & Detectar y alertar aglomeraciones (regla global) & Operador o administrador & Implementado \\
CU-06 & Visualizar el plano en vivo y el video por cámara & Operador o administrador & Implementado \\
CU-07 & Exportar reportes & Operador o administrador & Implementado \\
CU-08 & Gestionar zonas de interés & Operador & Implementado \\
CU-09 & Configurar reglas de aglomeración por zona & Operador & Implementado \\
CU-10 & Consultar el dashboard ejecutivo y el panel comercial & Operador o administrador & Implementado \\
CU-11 & Validar una fuente en modo laboratorio & Operador & Implementado \\
CU-12 & Consultar la auditoría del sistema & Operador o administrador & Implementado \\
CU-13 & Gestionar proyectos independientes & Operador o administrador & Implementado \\
CU-14 & Iniciar sesión según el rol & Operador o administrador & Implementado \\
CU-15 & Revisar la cámara en vivo antes de calibrar & Operador & Implementado \\
CU-16 & Configurar y recibir avisos de equipaje sin custodia & No aplica & Retirado \\
CU-17 & Configurar y recibir alertas de seguridad por correo & Operador o administrador; personal de seguridad & Implementado \\
CU-18 & Gestionar la bitácora de incidentes & Operador o administrador & Implementado \\
\end{longtable}
}

\begin{figure}[H]
\centering
\resizebox{\anchofigura}{!}{%
\begin{tikzpicture}[
  cu/.style={ellipse,draw=black,thick,minimum width=5.6cm,minimum height=1.05cm,
             text width=4.5cm,align=center,font=\footnotesize,inner sep=0pt},
  retirado/.style={cu,dashed},
  asoc/.style={black,thin}]
  \draw[thick] (-1.0,0.9) rectangle (13.6,-15.58);
  \node[font=\small\bfseries,anchor=north] at (6.3,0.8) {Sistema AeroTrack};
  \node[cu] (c01) at (2.0,-1.42) {CU-01\\Configurar plano y cámaras};
  \node[cu] (c02) at (2.0,-3.04) {CU-02\\Calibrar una cámara};
  \node[cu] (c03) at (2.0,-4.66) {CU-03\\Procesar fuente y seguir};
  \node[cu] (c08) at (2.0,-6.28) {CU-08\\Gestionar zonas de interés};
  \node[cu] (c09) at (2.0,-7.90) {CU-09\\Reglas por zona};
  \node[cu] (c11) at (2.0,-9.52) {CU-11\\Validar en laboratorio};
  \node[cu] (c15) at (2.0,-11.14) {CU-15\\Revisar cámara en vivo};
  \node[cu] (c05) at (10.6,-1.42) {CU-05\\Alertar aglomeraciones};
  \node[cu] (c06) at (10.6,-3.04) {CU-06\\Visualizar plano y video};
  \node[cu] (c07) at (10.6,-4.66) {CU-07\\Exportar reportes};
  \node[cu] (c10) at (10.6,-6.28) {CU-10\\Dashboard y panel comercial};
  \node[cu] (c12) at (10.6,-7.90) {CU-12\\Consultar auditoría};
  \node[cu] (c13) at (10.6,-9.52) {CU-13\\Gestionar proyectos};
  \node[cu] (c14) at (10.6,-11.14) {CU-14\\Iniciar sesión por rol};
  \node[cu] (c17) at (10.6,-12.76) {CU-17\\Alertas por correo};
  \node[cu] (c18) at (10.6,-14.38) {CU-18\\Bitácora de incidentes};
  \node[cu] (c04) at (2.0,-12.76) {CU-04\\Identidad y reidentificación};
  \node[retirado] (c16) at (4.6,-14.38) {CU-16 (retirado)\\Equipaje sin custodia};
  \begin{scope}[shift={(-2.90,-6.48)}]
    \draw[thick] (0,0.55) circle (0.22); \draw[thick] (0,0.33)--(0,-0.35);
    \draw[thick] (-0.35,0.1)--(0.35,0.1); \draw[thick] (0,-0.35)--(-0.3,-0.85); \draw[thick] (0,-0.35)--(0.3,-0.85);
    \node[font=\small,below,align=center] at (0,-0.9) {Operador};
  \end{scope}
  \begin{scope}[shift={(15.60,-8.10)}]
    \draw[thick] (0,0.55) circle (0.22); \draw[thick] (0,0.33)--(0,-0.35);
    \draw[thick] (-0.35,0.1)--(0.35,0.1); \draw[thick] (0,-0.35)--(-0.3,-0.85); \draw[thick] (0,-0.35)--(0.3,-0.85);
    \node[font=\small,below,align=center] at (0,-0.9) {Operador o\\administrador};
  \end{scope}
  \begin{scope}[shift={(-2.90,-13.36)}]
    \draw[thick] (0,0.55) circle (0.22); \draw[thick] (0,0.33)--(0,-0.35);
    \draw[thick] (-0.35,0.1)--(0.35,0.1); \draw[thick] (0,-0.35)--(-0.3,-0.85); \draw[thick] (0,-0.35)--(0.3,-0.85);
    \node[font=\small,below,align=center] at (0,-0.9) {Sistema\\(automático)};
  \end{scope}
  \coordinate (op) at (-2.5,-6.38);
  \coordinate (ad) at (15.2,-8.00);
  \coordinate (si) at (-2.5,-13.26);
  \foreach \id in {01,02,03,08,09,11,15} { \draw[asoc] (op) -- (c\id.west); }
  \foreach \id in {05,06,07,10,12,13,14,17,18} { \draw[asoc] (ad) -- (c\id.east); }
  \draw[asoc] (si) -- (c04.west);
  \draw[asoc,dashed] (si) -- (c16.west);
\end{tikzpicture}}
\caption{Diagrama de casos de uso del sistema AeroTrack. La línea discontinua señala el
módulo retirado del prototipo vigente.}
\label{fig:casos-uso}
\end{figure}

\subsubsection{Fichas de casos de uso}
\label{subsec:casos-uso-detalle}

\fichacu{CU-01}{Configurar el plano y las cámaras}
  {Operador.}
  {el servidor local está en ejecución, el usuario inició sesión con el rol de operador y hay un
   proyecto abierto.}
  {\item El operador elige uno de los planos del aeropuerto, crea un plano en blanco o importa
         una imagen, un PDF o un archivo DXF.
   \item Si importa una imagen o un PDF, recorta la región útil y el sistema la reduce a un dibujo
         esquemático de líneas por detección de bordes, en lugar de conservar la fotografía.
   \item Coloca cada cámara sobre el plano y define su alcance orientativo.
   \item Declara qué cámaras están relacionadas entre sí.}
  {el plano queda registrado con sus cámaras y disponible para la calibración (CU-02).}
  {}

\fichacu{CU-02}{Calibrar una cámara}
  {Operador.}
  {la cámara tiene una fuente de video válida.}
  {\item El operador marca un punto de referencia del suelo en el video y su correspondencia en
         el plano.
   \item Repite la marcación un mínimo de cuatro veces; se recomiendan de seis a ocho puntos
         repartidos por toda el área de interés.
   \item El sistema calcula la homografía y, con cinco o más puntos, comprueba cada referencia
         dejándola fuera del ajuste.}
  {la cámara queda habilitada para proyectar posiciones al plano compartido.}
  {el sistema rechaza la calibración si hay referencias repetidas, casi alineadas o
   geométricamente inconsistentes. La calibración es opcional: sin ella la cámara detecta y sigue
   personas en la imagen, pero no las ubica en el plano ni usa la posición para asociarlas entre
   cámaras.}

\fichacu{CU-03}{Procesar una fuente y hacer seguimiento}
  {Operador.}
  {la cámara tiene una fuente configurada (archivo de video, cámara USB o flujo de red) y su zona
   útil marcada; los pesos del detector YOLO y el modelo OSNet están instalados en el equipo.}
  {\item El operador inicia la sesión de monitoreo y elige el modo de rendimiento (automático,
         preciso, equilibrado o rápido).
   \item El sistema elige la variante de YOLO11 según el hardware del equipo y detecta personas en
         cada muestra de video.
   \item Descarta las detecciones fuera de la zona útil, las dudosas mucho más pequeñas que las
         personas seguras de la misma cámara y los objetos inmóviles de baja confianza.
   \item Sigue a cada persona con asociación en dos etapas sobre el punto de los pies, reforzada
         por la superposición de la caja predicha y una firma de color del torso, y le asigna un
         identificador local estable dentro de esa cámara.
   \item Si la cámara está calibrada, proyecta el punto de los pies al plano.}
  {cada persona detectada tiene un identificador local en esa cámara y, si la cámara está
   calibrada, una posición en el plano; las vistas confiables pasan al motor de identidad
   (CU-04).}
  {si falta el modelo OSNet, la sesión no se inicia. Si faltan los pesos de la variante de YOLO
   recomendada, el sistema usa la mejor variante disponible en el equipo y lo advierte.}

\fichacu{CU-04}{Asociar identidad entre cámaras y reidentificar personas}
  {Sistema (automático), a partir de la declaración de sincronización que hace el operador.}
  {sesión de monitoreo en curso. Para asociar personas entre cámaras: al menos dos cámaras
   relacionadas y la declaración expresa del operador de que los relojes de las fuentes están
   sincronizados. La calibración es opcional.}
  {\item Para cada persona seguida cuya vista es confiable (buena confianza, altura suficiente,
         lejos del borde y sin oclusión), el sistema calcula un vector de apariencia OSNet de 512
         dimensiones sobre el recorte de cuerpo entero; las partes tapadas por otra persona se
         rellenan con un gris neutro.
   \item Agrupa las vistas de cada tramo continuo de seguimiento y forma con ellas un prototipo
         de apariencia.
   \item La persona lleva un identificador provisional hasta reunir al menos cinco vistas durante
         dos segundos; entonces recibe un identificador confirmado, que es el único que cuenta en
         ocupación, flujo y alertas.
   \item Busca, entre las identidades de las cámaras relacionadas, una candidata con parecido de
         apariencia sobre el umbral y coherente en tiempo y, si hay calibración, en posición en el
         plano. Si hay una candidata inequívoca, reutiliza su identificador; si la mejor candidata es
         ambigua, marca la asociación como incierta en lugar de forzarla.
   \item Si la memoria de apariencia está activa, que es la configuración por defecto, compara
         cada identidad nueva con todas las personas guardadas dentro del plazo de retención y,
         si coincide sin ambigüedad, le devuelve su identificador anterior, aunque se trate de otra
         sesión u otro día.
   \item Al cerrar la sesión, reagrupa los tramos con la grabación completa a la vista y renumera
         a las personas contadas de 1 a N.}
  {cada observación queda con un identificador, provisional o confirmado, y un estado de
   asociación: local, estimada (une cámaras), reidentificada (recupera un tramo o una persona de
   la memoria) o incierta.}
  {sin sincronización declarada, el sistema no asocia personas entre cámaras en la sesión, pero
   la memoria de apariencia sigue comparando cada identidad nueva con las personas guardadas.
   La memoria conserva el vector en el disco del equipo durante 24 horas por defecto
   (configurable hasta 168), lo que constituye la reidentificación individual persistente que
   RNF-01 prohíbe (apartado~\ref{subsec:privacidad}).}

\fichacu{CU-05}{Detectar y alertar aglomeraciones (regla global)}
  {Operador o administrador.}
  {sesión de monitoreo en ejecución.}
  {\item Desde la vista de alertas, el usuario fija un radio en metros, una cantidad mínima de
         personas y un tiempo de permanencia, aplicables a todo el plano.
   \item El sistema evalúa en cada muestra las posiciones de las personas confirmadas.
   \item Marca una alerta cuando una concentración cumple los tres criterios de forma
         sostenida.}
  {la alerta queda registrada como incidente en la bitácora (CU-18) y se envía por correo si el
   aviso está configurado (CU-17).}
  {}

\fichacu{CU-06}{Visualizar el plano en vivo y el video por cámara}
  {Operador o administrador.}
  {sesión activa o grabación de una sesión anterior disponible.}
  {\item El usuario observa el plano con las personas, sus trayectorias recientes y las
         aglomeraciones resaltadas.
   \item Alterna al video real de una cámara, con la caja y el identificador de cada persona
         superpuestos.
   \item Puede revisar las sesiones grabadas desde la vista de videos y resultados.}
  {la consulta no modifica la configuración ni la sesión.}
  {}

\fichacu{CU-07}{Exportar reportes}
  {Operador o administrador.}
  {sesión con datos observados.}
  {\item El usuario elige el tipo de reporte: zonas de ocupación, comparación de sectores,
         ocupación por zona, flujo horario, trayectorias recientes o reporte comercial.
   \item Elige el formato: PDF, Excel, CSV o JSON.}
  {el sistema genera el archivo indicando siempre el origen de los datos: simulación sintética,
   videos de prueba o fuente real sin validación de precisión.}
  {el reporte de trayectorias contiene recorridos por identificador, por lo que hereda la
   calificación de dato personal de las grabaciones (Anexo~\ref{anx:marco}).}

\fichacu{CU-08}{Gestionar zonas de interés}
  {Operador.}
  {plano configurado.}
  {\item El operador dibuja polígonos, rectángulos o círculos sobre el plano.
   \item Clasifica cada área como zona de interés, cola, zona restringida, sala, pasillo, muro,
         puerta o zona comercial.
   \item Puede aceptar puntos sugeridos por el sistema según la ocupación observada.}
  {las zonas quedan registradas y disponibles para reglas y reportes por zona.}
  {}

\fichacu{CU-09}{Configurar reglas de aglomeración por zona}
  {Operador.}
  {al menos una zona dibujada.}
  {\item El operador activa una regla independiente para una zona, con su propia cantidad mínima
         de personas y su propio tiempo de permanencia.}
  {la zona genera alertas específicas según su uso; por ejemplo, una cola tolera más personas
   que una sala de espera.}
  {}

\fichacu{CU-10}{Consultar el dashboard ejecutivo y el panel comercial}
  {Operador o administrador.}
  {sesión con datos observados.}
  {\item El usuario revisa el ranking de zonas por ocupación acumulada, la evolución de las
         personas observadas y las zonas sugeridas para evaluación comercial.
   \item En el panel comercial consulta las zonas ordenadas por oportunidad comercial, el
         tráfico por acceso y las métricas de cada negocio (visitas, permanencias y regresos).}
  {la información se presenta agregada y el panel indica si procede de una fuente real o es
   sintética.}
  {aunque se presentan agregadas, las métricas por negocio se calculan a partir de las
   trayectorias por identificador, lo que constituye una finalidad distinta de la estimación de
   aglomeraciones (Sección~\ref{subsec:impacto-etico}).}

\fichacu{CU-11}{Validar una fuente en modo laboratorio}
  {Operador.}
  {cámara agregada al proyecto.}
  {\item El operador prueba la cámara de forma aislada: señal, resolución, frecuencia de cuadros,
         tiempo de procesamiento, calidad de la calibración y estabilidad del seguimiento, con
         identificadores temporales sobre el video.}
  {la cámara queda validada o rechazada antes de sumarse a una sesión real.}
  {}

\fichacu{CU-12}{Consultar la auditoría del sistema}
  {Operador o administrador.}
  {al menos una acción registrada.}
  {\item El usuario revisa el registro de acciones de la ejecución del servidor (guardar, iniciar,
         pausar, detener, abrir una vista en vivo o borrar la memoria de identidad, entre otras),
         con fecha, hora y detalle.}
  {queda disponible la trazabilidad operativa de la sesión, sin identidades reales.}
  {}

\fichacu{CU-13}{Gestionar proyectos independientes}
  {Operador; el administrador solo puede abrir proyectos existentes.}
  {servidor en ejecución con al menos un proyecto.}
  {\item El usuario ve la lista de proyectos desde la pantalla de inicio o el selector de la
         cabecera.
   \item El operador crea un proyecto nuevo, abre uno existente, lo renombra o lo elimina.}
  {el sistema cambia el proyecto activo; el plano, las cámaras, la calibración, los incidentes, las
   ventas y la memoria de apariencia de un proyecto nunca aparecen en otro.}
  {no se puede eliminar el único proyecto existente, y el sistema rechaza los cambios destinados
   a un proyecto distinto del abierto.}

\fichacu{CU-14}{Iniciar sesión según el rol}
  {Operador o administrador.}
  {existe al menos un usuario; si no existe ninguno, el primero que se registra queda como
   operador.}
  {\item El usuario ingresa su nombre y su contraseña.
   \item El sistema valida las credenciales contra el resumen criptográfico almacenado y abre una
         sesión con el rol correspondiente.}
  {el operador accede a la configuración y a todo lo demás; el administrador accede a todo
   menos a la configuración y puede ajustar los umbrales de alerta desde la vista de alertas.}
  {las credenciales incorrectas se rechazan sin indicar cuál de los dos datos falló.}

\fichacu{CU-15}{Revisar la cámara en vivo antes de calibrar}
  {Operador.}
  {la cámara tiene una fuente configurada y no hay una sesión de seguimiento en curso.}
  {\item El operador abre la vista en vivo de la cámara desde la etapa de cámaras del asistente.
   \item El sistema muestra la imagen en movimiento sin ejecutar detección; si la fuente es una
         grabación, permite pausar, reanudar y saltar a un instante.
   \item El operador puede obtener una imagen fija para marcar la zona útil.}
  {el operador confirma encuadre, luz y retardo; la apertura queda en el registro de auditoría y,
   con la imagen fija, la fuente queda comprobada con su resolución y su frecuencia de cuadros.}
  {si nadie solicita imágenes durante 20 segundos, el sistema libera la cámara.}

\fichacu{CU-16}{Configurar y recibir avisos de equipaje sin custodia}
  {no aplica en el prototipo vigente.}
  {no aplica.}
  {\item Este caso describía la detección de mochilas, bolsos y maletas inmóviles durante un
         tiempo configurable y la emisión de un aviso para revisión humana.}
  {no aplica.}
  {el módulo se retiró del prototipo al simplificarlo al esquema de seguimiento multicámara. Al
   abrir proyectos antiguos, su configuración se elimina. Se conserva el número para no alterar
   la numeración maestra. Los incidentes y los correos actuales se refieren solo a
   aglomeraciones.}

\fichacu{CU-17}{Configurar y recibir alertas de seguridad por correo}
  {Operador o administrador, que configura; personal de seguridad, que recibe.}
  {se cuenta con una cuenta de correo remitente y su contraseña de aplicación.}
  {\item Desde la vista de alertas, el usuario activa el aviso por correo e indica la cuenta
         remitente, su contraseña de aplicación y hasta veinte destinatarios.
   \item Puede enviar un correo de prueba.
   \item Ante cada incidente nuevo de aglomeración, el sistema envía un correo con el lugar, la
         cámara, el motivo del aviso y la advertencia de que requiere comprobación humana.}
  {la configuración queda guardada fuera de los archivos del proyecto, la contraseña no se
   muestra nunca en la interfaz y el cambio queda en el registro de auditoría.}
  {una misma zona no genera más de un correo cada 120 segundos y un aviso ya enviado no se
   reenvía; un fallo de envío queda registrado sin interrumpir la sesión.}

\fichacu{CU-18}{Gestionar la bitácora de incidentes}
  {Operador o administrador.}
  {existen incidentes de aglomeración registrados.}
  {\item El usuario consulta la lista de incidentes del proyecto activo, con su hora, zona,
         magnitud y estado.
   \item Marca cada incidente como revisado, falsa alarma o resuelto, o lo devuelve a
         pendiente.}
  {el incidente queda con su estado actualizado, lo que distingue las alertas atendidas de las
   pendientes.}
  {los incidentes se conservan aunque la sesión termine, y los de un proyecto nunca aparecen
   en otro.}

\subsubsection{Requisitos funcionales}
\label{subsec:rf-tabla}

{\small
\setlength{\LTleft}{0pt}\setlength{\LTright}{0pt}
\begin{longtable}{@{}L{1.3cm}L{6.6cm}L{4.2cm}L{1.9cm}@{}}
\caption{Requisitos funcionales y su estado verificado en el prototipo.}\label{tab:rf}\\
\toprule
\textbf{ID} & \textbf{Requisito} & \textbf{Estado} & \textbf{CU} \\
\midrule
\endfirsthead
\toprule
\textbf{ID} & \textbf{Requisito} & \textbf{Estado} & \textbf{CU} \\
\midrule
\endhead
\bottomrule
\endlastfoot
RF-01 & Configurar plano y cámaras, con recorte e importación como dibujo de líneas. & Implementado. & CU-01 \\
RF-02 & Calibrar por homografía con validación de consistencia geométrica. & Implementado. & CU-02 \\
RF-03 & Detectar personas con YOLO11, en la variante que corresponda al hardware, y seguirlas dentro de cada cámara con asociación en dos etapas sobre el punto de los pies. & Implementado; los pesos se instalan al preparar el sistema. & CU-03, CU-11 \\
RF-04 & Asociar identidad entre cámaras con el vector de apariencia OSNet y coherencia de tiempo y plano, solo entre cámaras relacionadas y con sincronización declarada, marcando \enquote{incierta} cuando corresponda, y reidentificar personas mediante la memoria de apariencia. & Implementado; su configuración por defecto incumple RNF-01. & CU-04 \\
RF-05 & Reglas de aglomeración globales configurables por radio, cantidad y permanencia. & Implementado. & CU-05 \\
RF-06 & Visualizar plano en vivo y video real por cámara con el identificador superpuesto. & Implementado. & CU-06 \\
RF-07 & Exportar reportes en PDF, Excel, CSV y JSON, por tipo (zonas, comparación de sectores, ocupación, flujo, trayectorias y comercial). & Implementado. & CU-07 \\
RF-08 & Dibujar y clasificar zonas de interés, incluyendo sugerencias automáticas por ocupación observada. & Implementado. & CU-08 \\
RF-09 & Configurar reglas de aglomeración independientes por zona. & Implementado. & CU-09 \\
RF-10 & Mostrar el dashboard ejecutivo con ranking de zonas y evolución temporal de personas observadas, y el panel comercial. & Implementado. & CU-10 \\
RF-11 & Validar una fuente de video en modo laboratorio antes de integrarla a la operación. & Implementado. & CU-11 \\
RF-12 & Registrar la auditoría de acciones del operador. & Implementado. & CU-12 \\
RF-13 & Gestionar proyectos independientes (crear, abrir, renombrar, eliminar), con aislamiento de datos entre ellos. & Implementado. & CU-13 \\
RF-14 & Autenticar usuarios con dos roles (operador y administrador) aplicados en el servidor, no solo ocultos en la interfaz. & Implementado. & CU-14 \\
RF-15 & Mostrar la vista en vivo de una cámara sin ejecutar detección, con controles de reproducción para grabaciones. & Implementado. & CU-15 \\
RF-16 & Detectar equipaje sin custodia por cámara, con tiempo de espera e intervalo de revisión configurables. & Retirado del prototipo vigente. & CU-16 \\
RF-17 & Enviar alertas de seguridad por correo electrónico, con control de repetición por incidente. & Implementado. & CU-17 \\
RF-18 & Registrar incidentes en una bitácora persistente con estado revisable (pendiente, revisado, falsa alarma, resuelto). & Implementado. & CU-18 \\
\end{longtable}
}

\subsection{Requerimientos de privacidad}
\label{subsec:privacidad}

La normativa de protección de datos exige que las garantías se incorporen desde la
concepción del tratamiento y que la configuración más protectora venga activada por defecto
(Sección~\ref{subsec:restricciones-legales}). El sistema recoge esa exigencia como requisitos no
funcionales, es decir, como condiciones que debe cumplir con independencia de su
funcionalidad.

El requisito rector es RNF-01, cuyo texto en la hoja maestra del equipo es: \enquote{No derivar
del video ningún rasgo que permita identificación biométrica o reidentificación individual
persistente}. \textbf{Al cierre de esta versión, RNF-01 es no conforme.} El sistema no construye
plantillas faciales, pero calcula para cada persona un vector de apariencia OSNet de cuerpo
entero y, con la memoria de apariencia activa, que es la configuración por defecto, lo guarda en
una base de datos local del equipo durante 24 horas (el plazo admite valores de 1 a 168 horas) y
lo usa para reconocer a la misma persona entre cámaras, entre sesiones y entre días. Eso es una
reidentificación individual persistente. La corrección está identificada y depende solo del
equipo: desactivar por defecto la memoria de apariencia, de modo que el vector viva
únicamente en memoria durante la sesión y se descarte con ella. Queda como acción pendiente.

De RNF-01 y del principio de proporcionalidad se derivan diez medidas de privacidad desde el
diseño. La Tabla~\ref{tab:medidas} las recoge con su estado en el prototipo.

{\small
\setlength{\LTleft}{0pt}\setlength{\LTright}{0pt}
\begin{longtable}{@{}L{1.2cm}L{5.8cm}L{7.4cm}@{}}
\caption{Medidas de privacidad desde el diseño: lo exigido y lo implementado.}\label{tab:medidas}\\
\toprule
\textbf{ID} & \textbf{Medida exigida} & \textbf{Estado en el prototipo} \\
\midrule
\endfirsthead
\toprule
\textbf{ID} & \textbf{Medida exigida} & \textbf{Estado en el prototipo} \\
\midrule
\endhead
\bottomrule
\endlastfoot
M-01 & Procesamiento en el borde, de modo que el fotograma no salga del perímetro del terminal. & Cumplida en el prototipo: todo el procesamiento ocurre en el equipo local y el servidor solo acepta conexiones del propio equipo. La localización del despliegue final sigue pendiente con LAP. \\
M-02 & No persistencia del fotograma, que se procesa en memoria y se descarta. & Cumplida para las cámaras en vivo. Los videos cargados para pruebas se conservan en el equipo. \\
M-03 & Salida agregada por zona, sin coordenadas individuales persistidas. & No cumplida: las grabaciones de cada sesión conservan, por muestra y por persona, el identificador, la caja y la posición en el plano, sin plazo de retención automático. \\
M-04 & Difuminado irreversible de rostros si algún fotograma debe conservarse. & No implementada: los videos de prueba se conservan sin difuminar. \\
M-05 & Plazo de retención definido, con purga automática. & Parcial: la memoria de apariencia se purga sola tras 24 horas; las grabaciones de sesión no tienen plazo automático y se borran a mano. \\
M-06 & Control de acceso por roles, con registro de auditoría. & Cumplida. \\
M-07 & Cifrado en tránsito y en reposo. & No implementada: la conexión con la consola es local y sin cifrar, y las bases de datos locales no están cifradas. \\
M-08 & Carteles informativos conforme al Anexo~\ref{anx:cartel}. & Pendiente: el modelo está propuesto, pero sus campos dependen de LAP. \\
M-09 & Documentación del linaje de los datos. & Parcial: el informe describe qué se calcula y dónde se guarda, sin un registro formal de linaje. \\
M-10 & Supervisión humana: el sistema alerta y no ejecuta acciones automáticas sobre personas. & Cumplida. \\
\end{longtable}
}

Donde un parámetro admita un ajuste más protector y otro menos protector (resolución de
captura, frecuencia de muestreo, granularidad de la salida, memoria de apariencia, plazo de
retención), el valor predeterminado debe ser el más protector. \textbf{Estas medidas describen el
sistema exigido y no el implementado}: una lista de medidas presentada al cliente como estado
alcanzado, cuando es estado deseado, es exactamente el tipo de afirmación que este informe
evita. A ello se suma una carencia de interfaz: el borrado inmediato de la memoria de identidad
existe como servicio interno del servidor, pero la consola no ofrece ningún botón para
ejecutarlo, de modo que un operador no puede atender una solicitud de supresión sin
asistencia técnica.

\subsection{Requerimientos no funcionales}
\label{subsec:rnf}

{\small
\setlength{\LTleft}{0pt}\setlength{\LTright}{0pt}
\begin{longtable}{@{}L{1.4cm}L{5.0cm}L{5.8cm}L{1.9cm}@{}}
\caption{Requisitos no funcionales y su estado verificado en el prototipo.}\label{tab:rnf}\\
\toprule
\textbf{ID} & \textbf{Requisito} & \textbf{Estado} & \textbf{CU} \\
\midrule
\endfirsthead
\toprule
\textbf{ID} & \textbf{Requisito} & \textbf{Estado} & \textbf{CU} \\
\midrule
\endhead
\bottomrule
\endlastfoot
RNF-01 & No derivar del video ningún rasgo que permita identificación biométrica o reidentificación individual persistente. & \textbf{No conforme.} La memoria de apariencia, activa por defecto, guarda el vector OSNet de cuerpo entero 24 horas (hasta 168) y reidentifica entre sesiones y días; además, las grabaciones persisten posiciones individuales (M-03). Acción pendiente: desactivar la memoria por defecto y fijar una retención automática para las grabaciones. & CU-03, CU-04, CU-06, CU-07, CU-10 \\
RNF-02 & Servidor accesible solo desde el propio equipo, con credencial de sesión en toda operación de escritura. & Implementado. & Transversal \\
RNF-03 & No persistir imágenes de forma permanente. & Cumplido para fuentes en vivo: los fotogramas se procesan en memoria y se descartan. Los videos cargados para pruebas se conservan en el equipo y su retención queda pendiente de definir. & CU-03, CU-06 \\
RNF-04 & Declarar expresamente cuándo el sistema opera en modo de demostración sintética. & Implementado. & CU-06, CU-07, CU-10 \\
RNF-05 & No presentar ninguna cifra de desempeño sobre datos públicos ni sobre grabaciones propias como desempeño esperado en el Aeropuerto Jorge Chávez sin validación en sitio. & Aplicado en este informe (Secciones~\ref{subsec:variables-restricciones} y~\ref{subsec:analisis-teas}). & CU-03, CU-11 \\
RNF-06 & Tiempo de inferencia compatible con la operación práctica. & No verificado como cumplido. Con YOLO11 nano y OSNet en CPU, analizar varias cámaras con una muestra cada 0,2 segundos va, según estimación del equipo, entre cinco y veinticinco veces más lento que el video; los modos de rendimiento espacian las muestras. Falta una medición con protocolo documentado. & CU-03, CU-11 \\
RNF-07 & Almacenar las contraseñas con resumen criptográfico y sal por usuario, nunca en texto plano ni expuestas por la interfaz de programación. & Implementado (PBKDF2 con SHA-256 y 240\,000 iteraciones). & CU-14 \\
RNF-08 & Los datos de un proyecto (plano, cámaras, incidentes, ventas, sesiones grabadas y memoria de apariencia) no deben ser accesibles desde otro proyecto. & Implementado y verificado con pruebas automatizadas. & CU-13, CU-18 \\
\end{longtable}
}

\subsection{Matriz de trazabilidad}
\label{subsec:trazabilidad}

La Tabla~\ref{tab:trazabilidad} relaciona cada requisito con los casos de uso que lo realizan.
Reproduce la matriz de la hoja maestra con una corrección: RNF-01 se vincula a los casos de uso
que calculan, conservan o explotan identificadores y trayectorias por persona (CU-03, CU-04,
CU-06, CU-07 y CU-10), en lugar de a la gestión de zonas y reglas.

\begin{table}[H]
\centering
\footnotesize
\setlength{\tabcolsep}{2.2pt}
\caption{Matriz de trazabilidad entre requisitos y casos de uso (columnas CU-01 a CU-18).}
\label{tab:trazabilidad}
\begin{tabular}{@{}lC{0.42cm}C{0.42cm}C{0.42cm}C{0.42cm}C{0.42cm}C{0.42cm}C{0.42cm}C{0.42cm}C{0.42cm}C{0.42cm}C{0.42cm}C{0.42cm}C{0.42cm}C{0.42cm}C{0.42cm}C{0.42cm}C{0.42cm}C{0.42cm}@{}}
\toprule
 & \multicolumn{18}{c}{\textbf{Caso de uso (CU-)}} \\
\cmidrule(l){2-19}
\textbf{Req.} & 01 & 02 & 03 & 04 & 05 & 06 & 07 & 08 & 09 & 10 & 11 & 12 & 13 & 14 & 15 & 16 & 17 & 18 \\
\midrule
RF-01 & X &  &  &  &  &  &  &  &  &  &  &  &  &  &  &  &  &  \\
RF-02 &  & X &  &  &  &  &  &  &  &  &  &  &  &  &  &  &  &  \\
RF-03 &  &  & X &  &  &  &  &  &  &  & X &  &  &  &  &  &  &  \\
RF-04 &  &  &  & X &  &  &  &  &  &  &  &  &  &  &  &  &  &  \\
RF-05 &  &  &  &  & X &  &  &  &  &  &  &  &  &  &  &  &  &  \\
RF-06 &  &  &  &  &  & X &  &  &  &  &  &  &  &  &  &  &  &  \\
RF-07 &  &  &  &  &  &  & X &  &  &  &  &  &  &  &  &  &  &  \\
RF-08 &  &  &  &  &  &  &  & X &  &  &  &  &  &  &  &  &  &  \\
RF-09 &  &  &  &  &  &  &  &  & X &  &  &  &  &  &  &  &  &  \\
RF-10 &  &  &  &  &  &  &  &  &  & X &  &  &  &  &  &  &  &  \\
RF-11 &  &  &  &  &  &  &  &  &  &  & X &  &  &  &  &  &  &  \\
RF-12 &  &  &  &  &  &  &  &  &  &  &  & X &  &  &  &  &  &  \\
RF-13 &  &  &  &  &  &  &  &  &  &  &  &  & X &  &  &  &  &  \\
RF-14 &  &  &  &  &  &  &  &  &  &  &  &  &  & X &  &  &  &  \\
RF-15 &  &  &  &  &  &  &  &  &  &  &  &  &  &  & X &  &  &  \\
RF-16 &  &  &  &  &  &  &  &  &  &  &  &  &  &  &  & X &  &  \\
RF-17 &  &  &  &  &  &  &  &  &  &  &  &  &  &  &  &  & X &  \\
RF-18 &  &  &  &  &  &  &  &  &  &  &  &  &  &  &  &  &  & X \\
\midrule
RNF-01 &  &  & X & X &  & X & X &  &  & X &  &  &  &  &  &  &  &  \\
RNF-02 & X & X & X & X & X & X & X & X & X & X & X & X & X & X & X & X & X & X \\
RNF-03 &  &  & X &  &  & X &  &  &  &  &  &  &  &  &  &  &  &  \\
RNF-04 &  &  &  &  &  & X & X &  &  & X &  &  &  &  &  &  &  &  \\
RNF-05 &  &  & X &  &  &  &  &  &  &  & X &  &  &  &  &  &  &  \\
RNF-06 &  &  & X &  &  &  &  &  &  &  & X &  &  &  &  &  &  &  \\
RNF-07 &  &  &  &  &  &  &  &  &  &  &  &  &  & X &  &  &  &  \\
RNF-08 &  &  &  &  &  &  &  &  &  &  &  &  & X &  &  &  &  & X \\
\bottomrule
\end{tabular}
\end{table}
